\section{Số học}
\begingroup\small
\begin{multicols}{2}
\subsection{Euclid mở rộng và nghịch đảo modulo}
Với $a,b>0$, thuật toán Euclid cho $\gcd(a,b)$ bằng cách lặp
$(a,b)\leftarrow(b,a\bmod b)$; độ phức tạp $O(\log\min(a,b))$.
Euclid mở rộng còn tìm $x,y$ sao cho
\[
ax+by=\gcd(a,b).
\]
Nếu $g=\gcd(a,m)=1$, hệ số $x$ thỏa $ax\equiv1\pmod m$,
nên $x\bmod m$ là nghịch đảo của $a$ modulo $m$.
Nghịch đảo tồn tại \emph{khi và chỉ khi} $g=1$; không cần $m$
nguyên tố. Mọi phép chia modulo $m$ phải nhân với nghịch đảo
của mẫu, và chỉ hợp lệ khi nghịch đảo tồn tại.
Nếu $m$ nguyên tố và $a\not\equiv0\pmod m$, cũng có thể dùng
$a^{m-2}\bmod m$ theo định lý Fermat.

Phương trình $ax\equiv b\pmod m$ có nghiệm khi và chỉ khi
$g=\gcd(a,m)$ chia hết $b$. Khi đó chia cả ba số cho $g$,
tìm một nghiệm modulo $m/g$, rồi có đúng $g$ nghiệm khác nhau
modulo $m$.

\subsection{Hệ phương trình đồng dư và CRT}
Xét $x\equiv r_1\pmod{m_1}$ và $x\equiv r_2\pmod{m_2}$,
với $m_1,m_2>0$. Đặt $g=\gcd(m_1,m_2)$. Hệ có nghiệm
\emph{khi và chỉ khi} $g\mid(r_2-r_1)$.
Viết $x=r_1+m_1t$ và chia đồng dư thứ hai cho $g$:
\[
\frac{m_1}{g}t\equiv\frac{r_2-r_1}{g}
\pmod{\frac{m_2}{g}}.
\]
Hai hệ số $m_1/g,m_2/g$ nguyên tố cùng nhau, nên tìm được
$t$ bằng Euclid mở rộng. Mọi nghiệm tạo một lớp modulo
$\operatorname{lcm}(m_1,m_2)$; ghép lần lượt nhiều phương trình.
Nếu các modulo đôi một nguyên tố cùng nhau, điều kiện tồn tại
luôn đúng và modulo kết quả là tích các modulo.
Kiểm tra tràn khi nhân modulo; chuẩn hóa phần dư về $[0,m)$,
kể cả khi dữ liệu ban đầu âm.

\paragraph{Mẫu C++17.}
\texttt{inverseMod} trả rỗng nếu không có nghịch đảo;
\texttt{mergeCRT} trả rỗng nếu hệ vô nghiệm, và ném lỗi nếu
modulo kết quả vượt \texttt{long long}. CRT cần modulo dương;
\texttt{inverseMod} chỉ xét $m\ge2$ và trả rỗng nếu $m\le1$.
\texttt{extendedGcd} nhận $a,b\ge0$, không đồng thời bằng $0$.
\cppsnippet{code/number_theory/crt.hpp}

\subsection{Số nguyên lớn trong C++ và Python}
Nếu kết quả vượt \texttt{long long} và không thể chỉ giữ modulo,
Python có kiểu \texttt{int} độ dài tùy ý; các phép
\texttt{+,-,*,//,\%}, \texttt{pow(a,b,m)} và
\texttt{math.comb(n,k)} dùng trực tiếp.
Trong C++, có thể dùng
\texttt{boost::multiprecision::cpp\_int}; ví dụ
\texttt{cpp\_int x = 1; x <<= 200;} tạo số $2^{200}$.
Kiểu này hỗ trợ cộng, trừ,
nhân, chia và nhập xuất như số nguyên thông thường, nhưng chi
phí thời gian và bộ nhớ lớn hơn kiểu máy. Phải kiểm tra quy định
ngôn ngữ/thư viện của kỳ thi trước khi dựa vào Boost hoặc Python.
Nếu chỉ cần tránh tràn \emph{trung gian}, dùng
\texttt{\_\_int128} rồi ép về kiểu nhỏ sau khi đã kiểm tra giới hạn.

\paragraph{BigNum tự cài bằng chuỗi (C++17).}
Lưu trị tuyệt đối dạng thập phân và một bit dấu; chuẩn hóa số $0$
thành \texttt{val="0", negative=false}. So sánh độ lớn theo số chữ số,
sau đó so từ điển. Cộng/trừ xử lý dấu rồi cộng/trừ trị tuyệt đối;
nhân từng chữ số, nhớ sang hàng kế tiếp. Cộng, trừ tốn $O(n)$;
nhân kiểu cơ bản tốn $O(nm)$ với hai toán hạng dài $n,m$ chữ số.
Mẫu hỗ trợ nhập, in, so sánh, \texttt{+}, \texttt{-}, \texttt{*};
không có phép chia. Khi được phép, \texttt{cpp\_int} thường tiện hơn.
\cppsnippet{code/number_theory/bignum.hpp}
\end{multicols}
\endgroup
